% !TeX program = XeLaTeX
% !TeX encoding = UTF-8
% !TeX spellcheck = pt_PT
% !TeX root = main
% IPBeja academic-work template — norms revised 2024.
% Compile with XeLaTeX and Biber (latexmk -xelatex main.tex).
\documentclass[12pt,a4paper,oneside,openright]{book}

% Language and fonts: Garamond is the institutional body-text requirement.
\usepackage[portuguese]{babel}
\usepackage{fontspec}
\setmainfont[Path=fonts/,Extension=.otf,UprightFont=texgyrepagella-regular,BoldFont=texgyrepagella-bold,ItalicFont=texgyrepagella-italic,BoldItalicFont=texgyrepagella-bolditalic]{texgyrepagella-regular}
\usepackage{microtype}

% Normative page layout: 2 cm left/right and 3 cm top/bottom.
\usepackage[a4paper,top=3cm,bottom=3cm,left=2cm,right=2cm]{geometry}
\usepackage{setspace}
\onehalfspacing
\setlength{\parindent}{1.27cm}
\setlength{\parskip}{0pt}

% Figures, tables, code and citations.
\usepackage{graphicx}
\usepackage{booktabs,tabularx}
\usepackage{listings}
\lstset{basicstyle=\ttfamily\small,breaklines=true,columns=fullflexible}
\usepackage{csquotes}
\usepackage[backend=biber,style=apa,sorting=nyt]{biblatex}
\DeclareLanguageMapping{portuguese}{portuguese-apa}
\addbibresource{main.bib}
\usepackage[hidelinks]{hyperref}

% Page numbers: lower right, 8 pt; Roman front matter, Arabic main matter.
\usepackage{fancyhdr}
\pagestyle{fancy}
\fancyhf{}
\fancyfoot[R]{\fontsize{8}{10}\selectfont\thepage}
\renewcommand{\headrulewidth}{0pt}
\renewcommand{\footrulewidth}{0pt}
\fancypagestyle{plain}{\fancyhf{}\fancyfoot[R]{\fontsize{8}{10}\selectfont\thepage}\renewcommand{\headrulewidth}{0pt}}

% Fill in the project-specific fields below. Keep cover free of decorative imagery.
\newcommand{\IPBInstitution}{Instituto Politécnico de Beja}
\newcommand{\IPBSchool}{Escola Superior de Tecnologia e Gestão}
\newcommand{\IPBProgramme}{[Ciclo de estudos]}
\newcommand{\IPBCourse}{[Unidade curricular]}
\newcommand{\IPBTitle}{[Título do trabalho]}
\newcommand{\IPBSubtitle}{[Subtítulo opcional]}
\newcommand{\IPBAuthor}{[Nome completo do/a estudante]}
\newcommand{\IPBSupervisor}{[Orientador/a, se aplicável]}
\newcommand{\IPBYear}{\the\year}

\begin{document}
\frontmatter

% Institutional cover: replace fields; place is always Beja.
\begin{titlepage}
  \centering
  {\Large\IPBInstitution\par}
  \vspace{1cm}
  {\large\IPBSchool\par}
  {\large\IPBProgramme\par}
  \vspace{1cm}
  {\large\IPBCourse\par}
  \vfill
  {\LARGE\bfseries\IPBTitle\par}
  \vspace{0.5cm}
  {\large\IPBSubtitle\par}
  \vfill
  {\large\IPBAuthor\par}
  \vspace{0.5cm}
  {\large Beja\par}
  {\large\IPBYear\par}
\end{titlepage}
\thispagestyle{empty}
\cleardoublepage

% Replace or remove optional front-matter sections as appropriate.
\chapter*{Folha de rosto}
\addcontentsline{toc}{chapter}{Folha de rosto}
\IPBTitle\par
\IPBSubtitle\par
\IPBAuthor\par
\IPBCourse\par
\IPBSchool\par
\IPBInstitution\par
\IPBSupervisor\par
\vspace{1em}
Trabalho realizado no âmbito da unidade curricular de \IPBCourse.
\cleardoublepage

\chapter*{Resumo}
\addcontentsline{toc}{chapter}{Resumo}
[Escrever um resumo com 100 a 250 palavras.]
\par\medskip\noindent\textbf{Palavras-chave:} [Três a seis palavras-chave; separadas por ponto e vírgula]
\cleardoublepage

\chapter*{Abstract}
\addcontentsline{toc}{chapter}{Abstract}
[Write an English abstract of 100 to 250 words.]
\par\medskip\noindent\textbf{Keywords:} [Three to six keywords; separated by semicolons]
\cleardoublepage

\tableofcontents
\listoffigures
\listoftables
% Add \lstlistoflistings only when the document contains listings.
\cleardoublepage

\mainmatter
% Chapters start on odd-numbered pages (openright); replace sample content.
\chapter{Introdução}
[Contextualização, problema, objetivos e organização do trabalho.]

\chapter{Desenvolvimento}
\section{Enquadramento}
[Desenvolver o tema e citar fontes no formato autor-data, por exemplo \parencite{ipbeja_normas_trabalhos_2024}.]
\section{Metodologia}
[Descrever métodos, materiais e procedimentos.]
\section{Resultados e discussão}
[Apresentar resultados, tabelas e figuras devidamente identificadas e citadas no texto.]

\chapter{Conclusão}
[Resumir conclusões, limitações e trabalho futuro.]

\backmatter
\printbibliography[heading=bibintoc,title={Referências Bibliográficas}]
% Optional: add glossary, appendices and annexes only when relevant.
\end{document}
